%%%%%%%%%%%%%%%%%%%%%%%%%%%%%%%%%%%%%%%%%%%%%%%%%%%%%%%%%%%%
%% Lecture 04
%%%%%%%%%%%%%%%%%%%%%%%%%%%%%%%%%%%%%%%%%%%%%%%%%%%%%%%%%%%%

\lecture
{04}
{Friday, 14 August 2026}
{Introduction to Probability}
{Dr. Nishant Chandgotia}

%%%%%%%%%%%%%%%%%%%%%%%%%%%%%%%%%%%%%%%%%%%%%%%%%%%%%%%%%%%%
\section{Sample Spaces, Random Variables and Their Distributions}

%%%%%%%%%%%%%%%%%%%%%%%%%%%%%%%%%%%%%%%%%%%%%%%%%%%%%%%%%%%%
\subsection{Sample Space}

The first object in probability is the \emph{sample space} $\Omega$,
which is the collection of all possible outcomes.

\begin{enumerate}

    \item \textbf{Simple coin toss}
    \(
    \qquad
    \Omega=\{H,T\}.
    \)

    \item \textbf{Arrival time of a bus}
    Depending on the context, we may take
    \(
    \Omega=[0,1]
    \)~~ 
    or~~
    \(
    \Omega=[0,\infty).
    \)

    \item \textbf{Infinite sequence of coin tosses}
    \(\qquad
    \Omega=\{H,T\}^{\mathbb{N}}.
    \)

\end{enumerate}

%%%%%%%%%%%%%%%%%%%%%%%%%%%%%%%%%%%%%%%%%%%%%%%%%%%%%%%%%%%%
\subsection{\texorpdfstring{$\sigma$-Algebra}{sigma-Algebra}}
\begin{enumerate}
    \item A $\sigma$-algebra $\mathcal{B}$ is the collection of events
whose probabilities we can measure.

For a simple coin toss,
\(\qquad
\mathcal{B}
=
\left\{
\varnothing,
\{H\},
\{T\},
\{H,T\}
\right\}.
\)

\item For a subset of $\mathbb{R}$, we are naturally led to the
\emph{Borel $\sigma$-algebra}.

\item For the infinite product space
\(\{H,T\}^{\mathbb{N}}\), consider a finite set
\(
F\subset\mathbb{N}
\)
and a function
\(
a:F\longrightarrow\{H,T\}.
\)

Define the corresponding cylinder set by

\[
[a]
=
\left\{
\omega\in\{H,T\}^{\mathbb{N}}
:
\omega_i=a(i)
\text{ for every }i\in F
\right\}.
\]

Cylinder sets form the basic building blocks for the product
$\sigma$-algebra.

    
\end{enumerate}
%%%%%%%%%%%%%%%%%%%%%%%%%%%%%%%%%%%%%%%%%%%%%%%%%%%%%%%%%%%%
\subsection{Probability Measure}

A probability measure is a measure
\(
\mathbb{P}:\mathcal{B}\longrightarrow[0,1]
\)
such that
\(
\mathbb{P}(\Omega)=1.
\)
\begin{enumerate}
    \item For a fair coin,
\( \qquad
\mathbb{P}(\{H\})
=
\mathbb{P}(\{T\})
=
\frac12.
\)

\item For the uniform probability measure on $[0,1]$,

\[
\mathbb{P}([a,b])=b-a,
\qquad
0\leq a\leq b\leq1.
\]

\item For independent coin tosses, if $[a]$ is a cylinder set determined
by $|F|$ coordinates, then

\[
\mathbb{P}([a])
=
2^{-|F|}.
\]

\item More generally, for a coin with
\(
\mathbb{P}(H)=p,
\qquad
\mathbb{P}(T)=1-p,
\)

we have

\[
\mathbb{P}([a])
=
p^{\#\text{Heads}}
(1-p)^{\#\text{Tails}}.
\]


Question: Why does this define a probability measure on
\(\{H,T\}^{\mathbb{N}}\)?
    
\end{enumerate}
%%%%%%%%%%%%%%%%%%%%%%%%%%%%%%%%%%%%%%%%%%%%%%%%%%%%%%%%%%%%
\subsection{Non-measurable Sets}

\begin{important}
Not every subset of
\(\{H,T\}^{\mathbb{N}}\)
is measurable.
\end{important}

The construction of non-measurable sets is related to the existence
of suitable equivalence relations and translation-invariant measures.

A useful analogy is the classical construction for $\mathbb{R}$.

%%%%%%%%%%%%%%%%%%%%%%%%%%%%%%%%%%%%%%%%%%%%%%%%%%%%%%%%%%%%

\subsubsection{Proof Idea on $\mathbb{R}$}

The general strategy is:

\begin{enumerate}

    \item Find a countable subgroup of $\mathbb{R}$.

    \item Find a Borel measure satisfying translation invariance:
    \[
    \mu(A+x)=\mu(A),
    \qquad
    \forall x\in\mathbb{R},
    \quad
    A\in\mathcal{B}.
    \]

\end{enumerate}

%%%%%%%%%%%%%%%%%%%%%%%%%%%%%%%%%%%%%%%%%%%%%%%%%%%%%%%%%%%%

\subsubsection{Analogue for the Infinite Coin-Toss Space}

Consider
\(\qquad
\left(\frac{\mathbb{Z}}{2\mathbb{Z}}\right)^{\mathbb{N}}.
\)

One can consider the set

\[
G
=
\left\{
x\in
\left(\frac{\mathbb{Z}}{2\mathbb{Z}}\right)^{\mathbb{N}}
:
x(i)=0
\text{ for all sufficiently large }i
\right\}.
\]

The idea is to check that $G$ is a countable subgroup and then
investigate translation invariance of the probability measure on
cylinder sets.

In particular, for a cylinder set $[a]$,
\[\qquad
\mathbb{P}([a])
=
2^{-|F|}.
\]

%\textbf{Todo:} Complete the construction and calculations.

%%%%%%%%%%%%%%%%%%%%%%%%%%%%%%%%%%%%%%%%%%%%%%%%%%%%%%%%%%%%
\section{Random Variables}

A random variable is a measurable function from the probability space
to another measurable space.

For example, for infinitely many coin tosses,
\(\qquad
f_n:
\{H,T\}^{\mathbb{N}}
\longrightarrow
[0,1],
\)
~~~ 
defined by
\[\qquad
f_n(\omega)
=
\frac{1}{n}
\sum_{i=1}^{n}X_i(\omega),
\]
where
\[
X_i(\omega)
=
\begin{cases}
1, & \omega_i=H,\\
0, & \omega_i=T.
\end{cases}
\]

We would like to understand the limit

\[
f(\omega)
=
\lim_{n\to\infty}f_n(\omega).
\]

%%%%%%%%%%%%%%%%%%%%%%%%%%%%%%%%%%%%%%%%%%%%%%%%%%%%%%%%%%%%
\subsection{A Random Variable Detecting the Existence of a Limit}

\begin{question}{Random Variable for Existence of Limit}{random-variable-question}
Given a probability space and a sequence of random variables
$(f_n)$, define a random variable which captures whether or not
\(
\lim_{n\to\infty}f_n(\omega)
\)
exists.
\end{question}

For a real-valued sequence, the limit exists precisely when
\(
\limsup_{n\to\infty}f_n(\omega)
=
\liminf_{n\to\infty}f_n(\omega)
\)
and both quantities are finite.

Define

\[
\widetilde f(\omega)
=
\begin{cases}
\displaystyle
\limsup_{n\to\infty}f_n(\omega)
-
\liminf_{n\to\infty}f_n(\omega),
&
\text{if both are finite},\\[2ex]
+\infty,
&
\text{otherwise}.
\end{cases}
\]

Then define

\[
f'(\omega)
=
\begin{cases}
1,
&
\widetilde f(\omega)=0,\\
0,
&
\widetilde f(\omega)\neq0.
\end{cases}
\]

Thus

\[
f'(\omega)
=
\mathbf{1}_{\{
\lim_{n\to\infty}f_n(\omega)
\text{ exists as a finite real number}
\}}.
\]

\begin{remark}

The precise formulation and notation may differ from the presentation
in Durrett.

\end{remark}

%%%%%%%%%%%%%%%%%%%%%%%%%%%%%%%%%%%%%%%%%%%%%%%%%%%%%%%%%%%%
\section{Distribution of a Random Variable}

Let~~~
\(
X:
(\Omega,\mathcal{B},\mathbb{P})
\longrightarrow
(S,\rho)
\)
be a random variable.

The distribution of $X$ is the push-forward of $\mathbb{P}$ under $X$:
\(
X_*\mathbb{P}.
\)

More explicitly, the distribution $\mu_X$ of $X$ is the measure on
$(S,\rho)$ defined by

\[
\mu_X(A)
=
\mathbb{P}\left(X^{-1}(A)\right),
\qquad
A\in\rho.
\]

It is also common to write

\[
\mathbb{P}(X\in A)
\coloneqq
\mathbb{P}\left(X^{-1}(A)\right).
\]

%%%%%%%%%%%%%%%%%%%%%%%%%%%%%%%%%%%%%%%%%%%%%%%%%%%%%%%%%%%%

Consider the random variables $X_1$ and $X_2$ arising from the
independent coin-toss example above.

Then $X_1$ and $X_2$ have the same distribution:~~~
\(
X_1\overset{d}{=}X_2.
\)

However, they are not the same random variable.

Indeed,
\(
X_1(\omega)
\neq
X_2(\omega)
\)
for some $\omega$.

Also,
\(
(X_1,X_2)
\)
is itself a random variable, taking values in the appropriate product
space.

\begin{remark}

Having the same distribution does not imply that two random variables
are equal.

In symbols,

\[
X=Y
\quad\Longrightarrow\quad
X\overset{d}{=}Y,
\]

but in general,

\[
X\overset{d}{=}Y
\quad\not\Longrightarrow\quad
X=Y.
\]

\end{remark}

\section{Some Useful Definitions:}
Let $(X_n)_{n\geq1}$ and $X$ be random variables defined on a
probability space $(\Omega,\mathcal{F},\mathbb{P})$.

%%%%%%%%%%%%%%%%%%%%%%%%%%%%%%%%%%%%%%%%%%%%%%%%%%%%%%%%%%%%
\subsection{Convergence in Probability}

\begin{definition}{Convergence in Probability}{conv-in-probability}
We say that $X_n$ \textbf{converges in probability} to $X$, denoted by
\[
    X_n \xrightarrow{\mathbb{P}} X,
\]
if for every $\varepsilon>0$,
\[
    \lim_{n\to\infty}
    \mathbb{P}\left(
        |X_n-X|>\varepsilon
    \right)
    =0.
\]

Equivalently, for every $\varepsilon>0$,
\[
    \mathbb{P}\left(
        |X_n-X|>\varepsilon
    \right)
    \longrightarrow 0.
\]
\end{definition}

%%%%%%%%%%%%%%%%%%%%%%%%%%%%%%%%%%%%%%%%%%%%%%%%%%%%%%%%%%%%
\subsection{Convergence in Distribution}

\begin{definition}{Convergence in Distribution}{conv-in-distribution}
We say that $X_n$ \textbf{converges in distribution} (or \emph{in law})
to $X$, denoted by
\[
    X_n \xrightarrow{d} X,
\]
if
\[
    \lim_{n\to\infty}F_{X_n}(x)
    =
    F_X(x)
\]
for every continuity point $x$ of the distribution function $F_X$ of
$X$.

Here,
\[
    F_{X_n}(x)
    =
    \mathbb{P}(X_n\leq x),
    \qquad
    F_X(x)
    =
    \mathbb{P}(X\leq x).
\]
\end{definition}
%%%%%%%%%%%%%%%%%%%%%%%%%%%%%%%%%%%%%%%%%%%%%%%%%%%%%%%%%%%%
\subsection{\texorpdfstring{Convergence in $L^1$}{Convergence in L1}}

\begin{definition}{\texorpdfstring{$L_1$}{L1} Convergence}{l1-conv}
    We say that $X_n$ \textbf{converges to $X$ in $L^1$}, denoted by
\[
    X_n \xrightarrow{L^1} X,
\]
if
\[
    \lim_{n\to\infty}
    \mathbb{E}\left[|X_n-X|\right]
    =0.
\]
Equivalently,
\[
    \|X_n-X\|_1
    =
    \mathbb{E}\left[|X_n-X|\right]
    \longrightarrow 0.
\]
\end{definition}

%%%%%%%%%%%%%%%%%%%%%%%%%%%%%%%%%%%%%%%%%%%%%%%%%%%%%%%%%%%%
\subsection{Almost Sure Convergence}

\begin{definition}{Almost Sure Convergence}{almost-sure-conv}
We say that $X_n$ \textbf{converges almost surely} to $X$, denoted by
\[
    X_n \xrightarrow{\mathrm{a.s.}} X,
\]
if
\[
    \mathbb{P}\left(
        \left\{
        \omega\in\Omega:
        \lim_{n\to\infty}X_n(\omega)=X(\omega)
        \right\}
    \right)
    =1.
\]

Equivalently,
\[
    \mathbb{P}\left(
        \lim_{n\to\infty}X_n=X
    \right)
    =1.
\]

That is, the sequence $X_n(\omega)$ converges to $X(\omega)$
for almost every $\omega\in\Omega$.
\end{definition}
%%%%%%%%%%%%%%%%%%%%%%%%%%%%%%%%%%%%%%%%%%%%%%%%%%%%%%%%%%%%
\section{Types of Random Variables}

Random variables can be classified according to their distributions.

\begin{enumerate}

    \item \textbf{Discrete random variables}

    \item \textbf{Continuous random variables}

    \begin{enumerate}

        \item Absolutely continuous random variables.

        \item Singular random variables.

    \end{enumerate}

\end{enumerate}

\begin{itemize}
    \item These classifications are determined by the distribution of the
random variable.

\item A discrete random variable has an \emph{atomic distribution}.
\end{itemize}

%%%%%%%%%%%%%%%%%%%%%%%%%%%%%%%%%%%%%%%%%%%%%%%%%%%%%%%%%%%%
\subsection{Examples of Discrete Random Variables}

\begin{enumerate}

    \item \textbf{Bernoulli random variable}
    Let
    \(
    p\in[0,1].
    \)~~We write
    \(
    X\sim\operatorname{Ber}(p)
    \)

    if
    \[
    \mathbb{P}(X=1)=p,
    \qquad
    \mathbb{P}(X=0)=1-p.
    \]

    This notation means that $X$ has the Bernoulli distribution with
    parameter $p$.

    \item \textbf{Binomial random variable}

    Let
    \(
    n\in\mathbb{N},
    \qquad
    p\in[0,1].
    \)
    We write
    \(
    X\sim\operatorname{Bin}(n,p)
    \)

    if
    \[
    \mathbb{P}(X=r)
    =
    \binom{n}{r}
    p^r(1-p)^{n-r},
    \qquad
    r=0,1,\ldots,n.
    \]

    \item \textbf{Poisson random variable}

    Let
    \(
    \lambda>0.
    \)
    We write~~ 
    \(
    X\sim\operatorname{Poi}(\lambda)
    \)

    if

    \[
    \mathbb{P}(X=k)
    =
    e^{-\lambda}
    \frac{\lambda^k}{k!},
    \qquad
    k=0,1,2,\ldots.
    \]

    \item \textbf{Geometric random variable}

    Let
    \(
    p\in(0,1).
    \)
    Under the convention that $X$ takes values in
    $\{1,2,\ldots\}$, we write
    \(
    X\sim\operatorname{Geo}(p)
    \)

    if

    \[
    \mathbb{P}(X=k)
    =
    p(1-p)^{k-1},
    \qquad
    k=1,2,\ldots.
    \]

\end{enumerate}

%%%%%%%%%%%%%%%%%%%%%%%%%%%%%%%%%%%%%%%%%%%%%%%%%%%%%%%%%%%%


\lectureend

\section{Quiz, 17 August 2026 }

\begin{quiz}

\begin{enumerate}
\item Convergence in probability implies convergence almost surely.

    \item Convergence in distribution implies convergence in probability.

    \item Convergence almost surely implies convergence in distribution.

    \item Let $X,Y,Z\sim\operatorname{Ber}(1/2)$ such that
    $X\perp Y$, $Y\perp Z$ and $Z\perp X$. Then

       $$ \mathbb{P}(X=1,Y=1,Z=1)=\frac{1}{8}.$$
\end{enumerate}

\end{quiz}


